\documentclass[10pt,a4paper]{amsart}
\usepackage{amsmath,amssymb,amsthm,graphicx}
\usepackage{verbatim}
\newenvironment{reference}{\comment}{\endcomment}
%\newenvironment{reference}{}{}
\newenvironment{slogan}{\comment}{\endcomment}
\newenvironment{history}{\comment}{\endcomment}
% For commutative diagrams we use Xy-pic
\usepackage[all]{xy}

% We use 2cell for 2-commutative diagrams.
\xyoption{2cell}
% Theorem environments.
%
\theoremstyle{plain}
\newtheorem{theorem}[subsection]{Theorem}
\newtheorem{proposition}[subsection]{Proposition}
\newtheorem{lemma}[subsection]{Lemma}

\theoremstyle{definition}
\newtheorem{definition}[subsection]{Definition}
\newtheorem{example}[subsection]{Example}
\newtheorem{exercise}[subsection]{Exercise}
\newtheorem{situation}[subsection]{Situation}

\theoremstyle{remark}
\newtheorem{remark}[subsection]{Remark}
\newtheorem{remarks}[subsection]{Remarks}

\numberwithin{equation}{subsection}

% Macros
%
\def\lim{\mathop{\mathrm{lim}}\nolimits}
\def\colim{\mathop{\mathrm{colim}}\nolimits}
\def\Spec{\mathop{\mathrm{Spec}}}
\def\Hom{\mathop{\mathrm{Hom}}\nolimits}
\def\Ext{\mathop{\mathrm{Ext}}\nolimits}
\def\SheafHom{\mathop{\mathcal{H}\!\mathit{om}}\nolimits}
\def\SheafExt{\mathop{\mathcal{E}\!\mathit{xt}}\nolimits}
\def\Sch{\mathit{Sch}}
\def\Mor{\mathop{\mathrm{Mor}}\nolimits}
\def\Ob{\mathop{\mathrm{Ob}}\nolimits}
\def\Sh{\mathop{\mathit{Sh}}\nolimits}
\def\NL{\mathop{N\!L}\nolimits}
\def\CH{\mathop{\mathrm{CH}}\nolimits}
\def\proetale{{pro\text{-}\acute{e}tale}}
\def\etale{{\acute{e}tale}}
\def\QCoh{\mathit{QCoh}}
\def\Ker{\mathop{\mathrm{Ker}}}
\def\Im{\mathop{\mathrm{Im}}}
\def\Coker{\mathop{\mathrm{Coker}}}
\def\Coim{\mathop{\mathrm{Coim}}}

% Boxtimes
%
\DeclareMathSymbol{\boxtimes}{\mathbin}{AMSa}{"02}

%
% Macros for moduli stacks/spaces
%
\def\QCohstack{\mathcal{QC}\!\mathit{oh}}
\def\Cohstack{\mathcal{C}\!\mathit{oh}}
\def\Spacesstack{\mathcal{S}\!\mathit{paces}}
\def\Quotfunctor{\mathrm{Quot}}
\def\Hilbfunctor{\mathrm{Hilb}}
\def\Curvesstack{\mathcal{C}\!\mathit{urves}}
\def\Polarizedstack{\mathcal{P}\!\mathit{olarized}}
\def\Complexesstack{\mathcal{C}\!\mathit{omplexes}}
% \Pic is the operator that assigns to X its picard group, usage \Pic(X)
% \Picardstack_{X/B} denotes the Picard stack of X over B
% \Picardfunctor_{X/B} denotes the Picard functor of X over B
\def\Pic{\mathop{\mathrm{Pic}}\nolimits}
\def\Picardstack{\mathcal{P}\!\mathit{ic}}
\def\Picardfunctor{\mathrm{Pic}}
\def\Deformationcategory{\mathcal{D}\!\mathit{ef}}
\usepackage{cleveref}
\crefname{originallemma}{original lemma}{original lemmas}
\Crefname{originallemma}{Original lemma}{Original lemmas}
\makeatletter
\newlabel{lemma-almost-exactness}{{lemma-almost-exactness}{source}{}{sourcepointer-1}{}}
\newlabel{lemma-almost-exactness@cref}{{[originallemma][0][]lemma-almost-exactness}{[1][0][]source}{}{}{}}
\makeatother
\crefname{originallemma}{original lemma}{original lemmas}
\Crefname{originallemma}{Original lemma}{Original lemmas}
\makeatletter
\newlabel{lemma-annihilated}{{lemma-annihilated}{source}{}{sourcepointer-2}{}}
\newlabel{lemma-annihilated@cref}{{[originallemma][0][]lemma-annihilated}{[1][0][]source}{}{}{}}
\makeatother
\crefname{originallemma}{original lemma}{original lemmas}
\Crefname{originallemma}{Original lemma}{Original lemmas}
\makeatletter
\newlabel{lemma-bound-q-and-d}{{lemma-bound-q-and-d}{source}{}{sourcepointer-3}{}}
\newlabel{lemma-bound-q-and-d@cref}{{[originallemma][0][]lemma-bound-q-and-d}{[1][0][]source}{}{}{}}
\makeatother
\crefname{originalreference}{original reference}{original references}
\Crefname{originalreference}{Original reference}{Original references}
\makeatletter
\newlabel{more-algebra-lemma-ses-near-projective}{{more-algebra-lemma-ses-near-projective}{source}{}{sourcepointer-4}{}}
\newlabel{more-algebra-lemma-ses-near-projective@cref}{{[originalreference][0][]more-algebra-lemma-ses-near-projective}{[1][0][]source}{}{}{}}
\makeatother
\title{Vanishing and comparison statements for the cohomology of a pullback of a uniformly annihilated module}
\author{The Stacks Project Authors}
\date{stacks-project:local-cohomology.tex:local-snapshot-20261012; GNU FDL 1.2 or later}
\begin{document}
\maketitle
\noindent\textit{Editorial note.} All delivered mathematical text is the author's own, copied from the pinned Stacks Project chapter source sources/stacks-local-cohomology/pinned-original-source.tex; no statement or proof step was rewritten, reordered or removed.\par\medskip
\noindent\textit{Editorial note.} Preamble provenance: the Stacks Project keeps one single shared preamble for every chapter; this card takes it directly from the staged sources/stacks-local-cohomology/preamble.tex as three verbatim spans (lines 8-12, 14-18 and 154-217), so no derived file was written and nothing was added or redefined. The wrapper class and the base packages are supplied by the assembler.\par\medskip
\noindent\textit{Editorial note.} Bibliography provenance: neither the counted statement nor its proof contains a citation (this lane rescans the whole span before writing the card), so the card carries no bibliography and no entry was reconstructed.\par\medskip
\noindent\textit{Editorial note.} Selection provenance: this card was opened under the file-based control contract (CONTROL.md), with the mechanical pre-card gate run on its statement and proof interval before the card was written.\par\medskip
\noindent\textit{Allowed standard prerequisites (editorial).}\begin{itemize}
\item The notion of an (A, n, c)-module and the invariants q\_0 = q(S) and d = d(S) of the graded situation. Set up in the author's section on a bit of uniformity and recalled in the statement itself. Fix the hypothesis on the module and the thresholds appearing in the five cases.
\item The behaviour of the class of (A, n, c)-modules under short exact sequences. Invoked inside the counted proof at source line 5176 and attributed to a lemma of another chapter (a cross-chapter pointer). Allows the proof to replace the given module by a resolution by free modules over the quotient and keep control of the constants.
\end{itemize}\medskip
\ifdefined\hypertarget\hypertarget{sourcepointer-1}{}\fi
\noindent\textit{Source reference.} Original lemma lemma-almost-exactness: sources/stacks-local-cohomology/pinned-original-source.tex, source lines 5190--5190. Navigation pointer only: the counted text refers to the Stacks item lemma-almost-exactness, which belongs to another chapter of the project and is not part of this card. The pointer carries the target's own label as its number because this chapter's staged snapshot does not contain the other chapter's tag table; its content is not reproduced here.\par\medskip
\ifdefined\hypertarget\hypertarget{sourcepointer-2}{}\fi
\noindent\textit{Source reference.} Original lemma lemma-annihilated: sources/stacks-local-cohomology/pinned-original-source.tex, source lines 5241--5241. Navigation pointer only: the counted text refers to the Stacks item lemma-annihilated, which belongs to another chapter of the project and is not part of this card. The pointer carries the target's own label as its number because this chapter's staged snapshot does not contain the other chapter's tag table; its content is not reproduced here.\par\medskip
\ifdefined\hypertarget\hypertarget{sourcepointer-3}{}\fi
\noindent\textit{Source reference.} Original lemma lemma-bound-q-and-d: sources/stacks-local-cohomology/pinned-original-source.tex, source lines 5200--5200. Navigation pointer only: the counted text refers to the Stacks item lemma-bound-q-and-d, which belongs to another chapter of the project and is not part of this card. The pointer carries the target's own label as its number because this chapter's staged snapshot does not contain the other chapter's tag table; its content is not reproduced here.\par\medskip
\ifdefined\hypertarget\hypertarget{sourcepointer-4}{}\fi
\noindent\textit{Source reference.} Original reference more-algebra-lemma-ses-near-projective: sources/stacks-local-cohomology/pinned-original-source.tex, source lines 5179--5179. Navigation pointer only: the counted text refers to the Stacks item more-algebra-lemma-ses-near-projective, which belongs to another chapter of the project and is not part of this card. The pointer carries the target's own label as its number because this chapter's staged snapshot does not contain the other chapter's tag table; its content is not reproduced here.\par\medskip
\begin{lemma}
\label{lemma-vanishing-coh-almost-projective}
With $q_0 = q(S)$ and $d = d(S)$ as above, suppose we have
integers $n \geq c \geq 0$, an $(A, n, c)$-module $M$,
an index $i \in \{0, 1, \ldots, d\}$, and an integer $q$.
Then we distinguish the following cases
\begin{enumerate}
\item In the case $i = d \geq 1$ and $q \geq q_0$ we have
$H^d(X, p^*M(q)) = 0$.
\item In the case $i = d - 1 \geq 1$ and $q \geq q_0$ we have
$H^{d - 1}(X, p^*M(q)) = 0$.
\item In the case $d - 1 > i > 0$ and $q \geq q_0 + (d - 1 - i)c$
the map
$H^i(X, p^*M(q)) \to H^i(X, p^*M(q - (d - 1 - i)c))$
is zero.
\item In the case $i = 0$, $d \in \{0, 1\}$, and $q \geq q_0$, there
is a surjection
$$
I^qM \longrightarrow H^0(X, p^*M(q))
$$
\item In the case $i = 0$, $d > 1$, and $q \geq q_0 + (d - 1)c$ the map
$$
H^0(X, p^*M(q)) \to H^0(X, p^*M(q - (d - 1)c))
$$
has image contained in the image of the canonical map
$I^{q - (d - 1)c}M \to H^0(X, p^*M(q - (d - 1)c))$.
\end{enumerate}
\end{lemma}

\begin{proof}
Let $M$ be an $(A, n, c)$-module. Choose a short exact sequence
$$
0 \to K \to (A/I^n)^{\oplus r} \to M \to 0
$$
We will use below that $K$ is an $(A, n, c)$-module, see More on Algebra,
Lemma \ref{more-algebra-lemma-ses-near-projective}.
Consider the corresponding exact sequence
$$
p^*K \to (\mathcal{O}_{Y_n})^{\oplus r} \to p^*M \to 0
$$
We split this into short exact sequences
$$
0 \to \mathcal{F} \to p^*K \to \mathcal{G} \to 0
\quad\text{and}\quad
0 \to \mathcal{G} \to (\mathcal{O}_{Y_n})^{\oplus r} \to p^*M \to 0
$$
By Lemma \ref{lemma-almost-exactness} the coherent module $\mathcal{F}$
is scheme theoretically supported on $Y_c$.

\medskip\noindent
Proof of (1). Assume $d > 0$. We have to prove
$H^d(X, p^*M(q)) = 0$ for $q \geq q_0$.
By the vanishing of the cohomology of twists of $\mathcal{G}$ in degrees $> d$
and the long exact cohomology sequence associated to the second
short exact sequence above, it suffices to prove that
$H^d(X, \mathcal{O}_{Y_n}(q)) = 0$.
This is true by Lemma \ref{lemma-bound-q-and-d}.

\medskip\noindent
Proof of (2). Assume $d > 1$. We have to prove
$H^{d - 1}(X, p^*M(q)) = 0$ for $q \geq q_0$.
Arguing as in the previous paragraph, we see that it suffices
to show that $H^d(X, \mathcal{G}(q)) = 0$. Using the first
short exact sequence and the vanishing of the cohomology
of twists of $\mathcal{F}$ in degrees $> d$ we see that it suffices
to show $H^d(X, p^*K(q))$ is zero which is
true by (1) and the fact that $K$ is an $(A, n, c)$-module (see above).

\medskip\noindent
Proof of (3). Let $0 < i < d - 1$ and assume the statement holds for $i + 1$
except in the case $i = d - 2$ we have statement (2).
Using the long exact sequence of cohomology associated to the second
short exact sequence above we find an injection
$$
H^i(X, p^*M(q - (d - 1 - i)c)) \subset
H^{i + 1}(X, \mathcal{G}(q - (d - 1 - i)c))
$$
as $q - (d - 1 - i)c \geq q_0$ gives the vanishing of
$H^i(X, \mathcal{O}_{Y_n}(q - (d - 1 - i)c))$
(see above). Thus it suffices to show that the map
$H^{i + 1}(X, \mathcal{G}(q)) \to H^{i + 1}(X, \mathcal{G}(q - (d - 1 - i)c))$
is zero. To study this, we consider the maps of exact sequences
$$
\xymatrix{
H^{i + 1}(X, p^*K(q)) \ar[r] \ar[d] &
H^{i + 1}(X, \mathcal{G}(q)) \ar[r] \ar[d] \ar@{..>}[ld] &
H^{i + 2}(X, \mathcal{F}(q)) \ar[d] \\
H^{i + 1}(X, p^*K(q - c)) \ar[r] \ar[d] &
H^{i + 1}(X, \mathcal{G}(q - c)) \ar[r] \ar[d] &
H^{i + 2}(X, \mathcal{F}(q - c)) \\
H^{i + 1}(X, p^*K(q - (d - 1 - i)c)) \ar[r] &
H^{i + 1}(X, \mathcal{G}(q - (d - 1 - i)c))
}
$$
Since $\mathcal{F}$ is scheme theoretically supported on $Y_c$
we see that the canonical map
$\mathcal{G}(q) \to \mathcal{G}(q - c)$ factors through
$p^*K(q - c)$ by Lemma \ref{lemma-annihilated}.
This gives the dotted arrow in the diagram. (In fact, for the proof it
suffices to observe that the vertical arrow on the extreme right is
zero in order to get the dotted arrow as a map of sets.)
Thus it suffices to show that
$H^{i + 1}(X, p^*K(q - c)) \to
H^{i + 1}(X, p^*K(q - (d - 1 - i)c))$
is zero. If $i = d - 2$, then the source of this arrow
is zero by (2) as $q - c \geq q_0$ and $K$ is an $(A, n, c)$-module.
If $i < d - 2$, then as $K$ is an $(A, n, c)$-module, we get from the
induction hypothesis that the map is indeed zero
since $q - c - (q - (d - 1 - i)c) = (d - 2 - i)c = (d - 1 - (i + 1))c$
and since $q - c \geq q_0 + (d - 1 - (i + 1))c$.
In this way we conclude the proof of (3).

\medskip\noindent
Proof of (4). Assume $d \in \{0, 1\}$ and $q \geq q_0$.
Then the first short exact sequence gives a surjection
$H^1(X, p^*K(q)) \to H^1(X, \mathcal{G}(q))$
and the source of this arrow is zero by case (1). Hence
for all $q \in \mathbf{Z}$ we see that the map
$$
H^0(X, (\mathcal{O}_{Y_n})^{\oplus r}(q))
\longrightarrow
H^0(X, p^*M(q))
$$
is surjective. For $q \geq q_0$ the
source is equal to $(I^q/I^{q + n})^{\oplus r}$ by
Lemma \ref{lemma-bound-q-and-d} and this easily proves the statement.

\medskip\noindent
Proof of (5). Assume $d > 1$. Arguing as in the proof of (4) we see that
it suffices to show that the image of
$$
H^0(X, p^*M(q))
\longrightarrow
H^0(X, p^*M(q - (d - 1)c))
$$
is contained in the image of
$$
H^0(X, (\mathcal{O}_{Y_n})^{\oplus r}(q - (d - 1)c))
\longrightarrow
H^0(X, p^*M(q - (d - 1)c))
$$
To show the inclusion above, it suffices to show that for
$\sigma \in H^0(X, p^*M(q))$ with boundary
$\xi \in H^1(X, \mathcal{G}(q))$ the image of $\xi$ in
$H^1(X, \mathcal{G}(q - (d - 1)c))$ is zero. This follows by the
exact same arguments as in the proof of (3).
\end{proof}

\end{document}
